%% ============================================================
%% ABSTRACT
%% Physics-Informed Game-Theoretic Defense of Pipeline Infrastructure
%% ============================================================

\chapter*{Abstract}
\addcontentsline{toc}{chapter}{Abstract}
\thispagestyle{plain}

\noindent\textbf{Title:}
\textit{Physics-Informed Game-Theoretic Defense of Pipeline Infrastructure:
Integrating Weld Joint Vulnerability into Bayesian Stackelberg Attacker--Defender
Models}

\medskip
\noindent\textbf{Author:} Babak Pirzadi \hfill
\textbf{Programme:} STRATEGOS MSc

\vspace{0.8em}
\noindent\rule{\textwidth}{0.4pt}
\vspace{0.6em}

%% ---- BODY (target: 280–320 words) --------------------------

Contemporary security game models for critical infrastructure treat attacker
payoffs as abstract constants, decoupled from the physical properties of
the assets they represent.  For high-consequence systems such as gas
transmission pipelines, this abstraction is a strategic liability: an
adversary who understands fracture mechanics and weld seam metallurgy will
target assets whose physics makes them disproportionately vulnerable, while
a defender applying game-theoretic resource allocation ignores the same
information.  This thesis closes that gap by constructing a four-layer,
physics-informed attacker--defender framework validated against
U.S.\ \ac{PHMSA} incident data spanning 2010 to 2024 (\num{1996} records,
216 material failure events).

The foundation (\emph{Zone C}) is a Monte Carlo failure probability engine
grounded in the \ac{BS7910}:2019 Level~2 \ac{FAD} and \ac{IIW} fatigue
classification system.  Weld seam--specific \acp{SCF} — calibrated from
\ac{PHMSA} fleet statistics — drive per-segment $\Pf$ values across a
20-node, 22-segment synthetic Gulf Coast transmission network.  The
resulting $\Pf$ range of $[\num{0.29},\; \num{0.93}]$ faithfully
reproduces the empirical failure-rate hierarchy: seamless ($\Pf \approx
\num{0.53}$) to fillet-welded legacy pipe ($\Pf \approx \num{0.88}$).
These physics-derived values are then used directly as attacker utilities
in a Bayesian Stackelberg Security Game with three adversary types
(\textsc{Strategic}, prior $= \num{0.50}$; \textsc{Opportunistic},
$\num{0.30}$; \textsc{State Actor}, $\num{0.20}$), solved via linear
programming enumeration.  At a defender budget of \SI{30}{\percent}, the
physics-informed \ac{SSE} allocation achieves \SI{+52.3}{\percent}
coverage effectiveness over the zero-budget baseline and a
\SI{17.0}{\percent} reduction in total expected network loss versus uniform
resource distribution.

The \emph{adversarial NDE threat layer} (\emph{Zone A}) examines whether
a sophisticated adversary could suppress the physical vulnerability signal
at the point of inspection.  A \num{32}-feature weld defect classifier
(\textsc{WeldDefectMLP}: $32{\to}128{\to}64{\to}4$, trained to
\SI{99.7}{\percent} clean accuracy) is subjected to white-box
$L_\infty$-bounded attacks: \ac{FGSM} achieves an \ac{ASR} of
\SI{18.1}{\percent} at $\eps = \num{0.30}$, \ac{BIM} \SI{20.1}{\percent},
and \ac{PGD} \SI{9.4}{\percent}~\cite{Goodfellow2015,Kurakin2017,Madry2018}.
The gap between BIM and PGD highlights that local gradient structure, not
random initialisation, dominates in this feature domain — a finding with
direct implications for NDE audit design.

The unified framework is delivered as a three-tab interactive \emph{Plotly
Dash} dashboard with click-to-inspect FAD assessment and segment-level
adversarial threat scoring, validated by a 310-test suite with full
coverage across all public interfaces.  The methodology is directly
extensible to Mediterranean transmission assets, including TAP and
SNAM~Rete~Gas corridors, using the same calibration pipeline with
locally-sourced incident data.

\vspace{0.6em}
\noindent\rule{\textwidth}{0.4pt}

\vspace{0.5em}
\noindent\textbf{Keywords:}
pipeline security, Bayesian Stackelberg game, BS~7910 fracture assessment,
IIW fatigue, adversarial machine learning, non-destructive examination,
critical infrastructure protection, PHMSA calibration, FGSM, PGD,
Monte Carlo failure probability
